% Build the layer, Chapter 7, second build: a claim with a bond behind it.
% Inserted after the first build (additions/bonded-build-the-ledger.tex),
% at the end of \S\ref{sec:conservation}. It builds the safety property of
% the session lifecycle of \S\ref{sec:tlaplus}, which that section says no
% model checker has run.
%
% No code or output in this file is typed. \pdexcerpt, \pdrun and
% \pdrunlines set files under excerpts/, which the kernel's
% tools/excerpts.mjs cuts from its source and its recorded runs (see
% figures/pd-listing.tex). The figure draws two of those recorded tables as
% they were printed. To change a listing or the figure, change the kernel
% and cut again. The prose names functions, never line numbers, so that a
% listing which moves in its file does not make a sentence wrong.
%
% \pdlistingroom{n} before a listing of n lines keeps it on one page. When a
% listing changes length, change the number with it.
\subsection{Build the layer: a claim with a bond behind it}
\label{sec:bonded-stake}

\noindent\emph{Express lane: the lifecycle's rule, that no work goes on
without units in escrow, is built twice. The first harbor is made of two
parts this book has already checked, with a careful desk between them. In
it a worker takes back the bond of a claim and keeps the claim. The second
harbor gives the claim and posts the bond in one decision. A walk
takes it through all 7,639 states of two workers and two paths within
small bounds, and a solver takes one step of an induction for any numbers. One of its guards closes a
gap that Chapter~1 left open on purpose.}

\S\ref{sec:tlaplus} states the session lifecycle in words, and says that no
model checker has run it. Its safety property is that every active session
carries positive escrow: an agent may not begin without posting it, and
nothing lowers it while the session is at work. \pdkernelname{} now has
both halves of that sentence. Chapter~1 gave it claims
(\weld{swk:sec:swk-build}{lease}), and a worker at work on a path is one
that holds a claim on it. The last section gave it funds and bonds. This
section asks what it takes to make the sentence true of the two together.

It takes more than having both. \S\ref{sec:conservation} remarks that
positive escrow and conservation are independent, and the same holds one
level up. The rules for claims and the rules for funds can each be right,
and a sentence about both can still be false, because it is a rule that
neither part owns.

\subsubsection{Naive: two right parts and a desk}

The obvious way to build this is out of what is already there. Call the
thing a worker talks to the \emph{desk}. For a claim the desk asks
Chapter~1's rules, and for a bond it asks the keeper of the last section.
Both answer through the decision core. The desk decides nothing about
either, and only keeps the two in step
(Listing~\ref{alg:bonded-stake-apart}).

\pdlistingroom{19}
\pdexcerpt[caption={The desk gives a claim. It asks for the claim, and then
posts the bond.},label={alg:bonded-stake-apart}]{ch07b-apart}

This desk is careful. A claim that is renewed keeps the bond it has, so a
worker is not charged again for staying. A grant is kept only if its bond
posts, so a project with an empty wallet gets no claims. Every rule the
desk leans on is one that an earlier section has checked.

The bond it posts is the harbor's \emph{price}, the units a claim must put
up. How to set a price is the question of \S\ref{sec:pricing}, and this
section does not take it up. In the runs below it is 6, the bond of the
chapter's repair.

\subsubsection{Break: a claim with nothing behind it}

Wren claims a path, and then asks the keeper for the bond back
(Listing~\ref{alg:bonded-stake-run-apart}).

\pdlistingroom{17}
\pdrun[caption={The break. Wren holds two claims, and escrow holds
nothing.},label={alg:bonded-stake-run-apart}]{free-hold-apart}

Step 3 is no mistake of the keeper's. Bond 1 is held and has not been
settled, and that is all the guard of Listing~\ref{alg:bonded-build-held}
asks. The keeper does not know that a claim stands on that bond, and it has
no way to find out. So the wallet is whole again and the claim is still
Wren's. Steps 4 and 5 do it a second time with the same units, and at
step 6 Heron is refused a path that costs Wren nothing to hold. A wallet
that covers one bond now stands behind as many claims as Wren cares to
make. \S\ref{sec:economic} asked the opposite, that each of a hundred
identities should need a bond of its own.

\subsubsection{Break again: a change after the bond has gone back}

The desk's own care opens a second hole. It refunds the bond when the
claim is released, which is right. Then Wren changes the path once more
(Listing~\ref{alg:bonded-stake-run-after}).

\pdlistingroom{9}
\pdrunlines[caption={The second break. The release gives the bond back,
and the change after it is admitted.},label={alg:bonded-stake-run-after}]{after-release-apart}{1,10-17}{1 and 10\textendash 17}

Step 5 is no mistake of the fence's either. Chapter~1's fence asks whether
a token is the newest one for its path, in the hands of the worker it was
given to (\weld{swk:sec:swk-build}{fence}). It does not ask whether the
grant was released. That was right in Chapter~1: until someone else
claims the path, a late write from its last holder hurts no claim. On this harbor it is a
change with no bond behind it. If it breaks something,
Theorem~\ref{thm:cleanup-bound} has nothing to pay for the cleanup with.

\subsubsection{Fix: one decision for both}

Neither part can be taught the missing rule, because the rule is about
both. \pdkernelname{} gives it a place of its own in the decision core
(\weld{prereq:sec:prereq-build}{decision core}). The core gets a record
that ties one grant to one bond, and calls it a \emph{stake}
(Listing~\ref{alg:bonded-stake-stake}).

\pdlistingroom{7}
\pdexcerpt[caption={A stake: one grant's token, and the bond behind
it.},label={alg:bonded-stake-stake}]{ch07b-stake}

A claim is now decided in two steps inside one decision. Chapter~1's rules
say whether the claim is given and what token it gets. Then \texttt{back}
puts a bond behind that grant (Listing~\ref{alg:bonded-stake-back}). The
answer carries the grant and its bond together, and the kernel stores both
before it returns. If the wallet cannot cover the bond,
\texttt{post} refuses, the question mark after it hands that refusal back,
and no grant is given.

\pdlistingroom{27}
\pdexcerpt[caption={The decision for a claim, once Chapter~1's rules have
given the grant. \texttt{own} is its
guard.},label={alg:bonded-stake-back}]{ch07b-back}

\texttt{own} looks slight. The kernel shows the core the stake of the
path's newest grant, and \texttt{own} asks whether that grant is the one
just given. For a renewal it is, since a renewal keeps its token, and the
bond is kept. For any other grant a new bond is posted. The section ends
with what happens without this guard.

A release runs the same way round. Chapter~1's rules release the grant,
and \texttt{free} gives back the bond behind it
(Listing~\ref{alg:bonded-stake-free}). It does so through the
\texttt{refund} of the last section, so every guard there still applies.

\pdlistingroom{19}
\pdexcerpt[caption={The decision for a release. The bond comes back
once.},label={alg:bonded-stake-free}]{ch07b-free}

That leaves the two holes. The first was a refund asked for by itself.
\texttt{loose} stands in front of every refund and every slash
(Listing~\ref{alg:bonded-stake-loose}).

\pdlistingroom{19}
\pdexcerpt[caption={The guard on a refund or a slash that is asked for by
itself.},label={alg:bonded-stake-loose}]{ch07b-loose}

The second was a change after the bond had gone back. The fence is left as
Chapter~1 wrote it, and \texttt{at\_stake} asks one more question after it
(Listing~\ref{alg:bonded-stake-at-stake}). It does not ask whether the
grant was released. It asks for the bond, which is the thing that matters,
and so it also refuses a change whose bond has gone for any other reason.

\pdlistingroom{21}
\pdexcerpt[caption={The guard on a change. A harbor whose price is nothing
asks for no bond, and is
Chapter~1's.},label={alg:bonded-stake-at-stake}]{ch07b-at-stake}

Wren's second run now stops at step 5
(Listing~\ref{alg:bonded-stake-run-after-core}).

\pdlistingroom{6}
\pdrunlines[caption={The fix. The change after the release is
refused.},label={alg:bonded-stake-run-after-core}]{after-release-core}{1,13-17}{1 and 13\textendash 17}

\keyidea[Give a rule about two things one owner]{Two parts that are each
right can be wrong together, when a rule mentions both and neither holds
it. Care in the code that joins them does not mend this, because each part
still answers by its own rule to whoever asks it directly. The rule needs
one place that decides for both, and every request has to pass through
it.}

\Needspace{12\baselineskip}
\subsubsection{Verify: every state, and then any numbers}

Wren's first run again, from step 3
(Listing~\ref{alg:bonded-stake-run-core}).

\pdlistingroom{15}
\pdrunlines[caption={The fix. The refund is refused, and a second claim
needs a second bond.},label={alg:bonded-stake-run-core}]{free-hold-core}{1,8-21}{1 and 8\textendash 21}

\paragraph{A claim that is left.}
One case is new. A lease can run out with its claim never released, and
another worker can then take the path
(Listing~\ref{alg:bonded-stake-run-left}).

\pdlistingroom{20}
\pdrun[caption={A claim left to run out. Its bond stays where it
is.},label={alg:bonded-stake-run-left}]{walked-away-core}

Heron's claim posts a bond of its own at step 4. Wren's token is now an old
one, so the fence refuses the release, and \texttt{loose} refuses the
refund. Bond 1 is held, and nothing in this build can move it. That is on
purpose. \S\ref{sec:pricing:settlement} says that a missed heartbeat
decides no slash, and that a bond in doubt is held for someone with the
authority to rule. A harbor in which step 6 went through would have made
Wren that judge. Had nobody else claimed the path, Wren could still have
released its claim and had the bond back, because a lease that runs out is
not a breach.

\paragraph{Every state.}
The walk of the last section is run again on the larger harbor, with two
workers and two paths at a price of 2 units. It tops up at most 4 units
and posts at most 3 bonds, and a path is given at most 2 grants. From each
state it tries every move the bounds allow, and most of them should be
refused. Either worker
claims a path the other holds. Either releases or changes a path under
every token ever given for it, and under one that never was. Every bond is
asked for by itself. Time is a move too, and it runs every lease out. Each
answer is held to a second set of rules: the lifecycle's sentence and
Chapter~1's, written out on plain values in a file that shares no code
with the core.

\pdlistingroom{21}
\pdrun[caption={Every state each harbor can reach, held to the lifecycle's
own rules.},label={alg:bonded-stake-states}]{stakes}

The kernel's harbor reaches 7,639 states and takes none of the moves that
the rules refuse. In every state a claim that has not been released has a
bond held behind it, and no bond stands behind two claims. In 480 of them
it holds a bond that waits, as bond 1 does above.

The harbor made of two parts is walked along the same road, and takes
25,272 moves that the rules refuse. It comes to 17,120 states in which a
claim has no bond, and it admits 7,032 changes under a claim with none.
The listing's last line matters as much as that row. Through all of this
the funds add up and no bond is shared, because each part keeps its own
rule in every state.

Figure~\ref{fig:bonded-pairs} draws what the two walks found.

% Brief (grammar: a small table of the possible combinations, with a mark
%   on each one that occurred; atlas VI/fig:bonded-pairs).
% Reader question: which combinations of a claim's status and its bond's
% status occur, in a harbor where a claim requires a bond?
% Claim: in the kernel's harbor four combinations occur, and in each the
% bond is held while the claim is not released. A harbor that decides
% claims and funds separately reaches six more with its first wrong
% operation, each a claim with no held bond, and it admits writes after a
% bond is refunded.
% Evidence: every state each harbor reached in the exhaustive test
% [internal, harbor stakes --each; sessions stakes-each.txt and
% stakes-each-separate.txt, drawn as printed, and
% tests/harbor-research/test_bonded_stake_numbers.py compares the counts
% in the caption with them].
% Must distinguish: a combination the reference model allows from one it
% does not; a claim under which writes were admitted from one under which
% none were.
% Grammar chosen: the claim's status in rows and the bond's in columns, so
% that the allowed combinations are a fixed set of cells, which are tinted,
% and a mark outside the tint is visibly a fault.
% Grammar rejected: a state diagram of one claim, which shows the order of
% operations and hides which combinations occur; bars of the counts, whose
% sizes depend on the bounds and mean nothing.
% Five-second test: in which table is there a mark outside the tint?
\begin{figure}[H]
\centering
\pdfigurehue{pdcobalt}%
\begin{tikzpicture}[pd figure,x=1.4cm,y=.6cm,
    allowed/.style={fill=pdfocus!11},
    occurs/.style={circle,fill=pdink,inner sep=0pt,minimum size=4.4pt},
    unbonded/.style={circle,fill=pderror,inner sep=0pt,minimum size=4.4pt},
    written/.style={draw=pdink,line width=.5pt},
    written unbonded/.style={draw=pderror,line width=.8pt}]
  % The tables are the recorded runs, read as printed: the first line is
  % the command, and every comment line starts with '#'.
  \pgfplotstableread[skip first n=1]{excerpts/stakes-each.txt}\pdpairsreal
  \pgfplotstableread[skip first n=1]{excerpts/stakes-each-separate.txt}\pdpairsseparate
  % Counts computed from the tables as they are drawn.
  \ifdefined\pdpairsin\else \newcount\pdpairsin \newcount\pdpairsout \fi
  % One panel. A claim's status is a row, with live at the top. A bond's
  % status is a column.
  \newcommand{\pdpairspanel}[1]{%
    \global\pdpairsin=0 \global\pdpairsout=0
    % The combinations the model allows: a held bond with a claim that is
    % not released, and a refunded bond with one that is.
    \foreach \y in {3,2,0}{ \fill[allowed] (.5,\y-.5) rectangle (1.5,\y+.5); }
    \fill[allowed] (1.5,.5) rectangle (2.5,1.5);
    \foreach \x in {-.5,.5,1.5,2.5,3.5}{ \draw[pdinkmuted!35,line width=.3pt] (\x,-.5) -- (\x,3.5); }
    \foreach \y in {-.5,.5,1.5,2.5,3.5}{ \draw[pdinkmuted!35,line width=.3pt] (-.5,\y) -- (3.5,\y); }
    \pgfplotstablegetrowsof{#1}
    \pgfmathtruncatemacro{\pdlastpair}{\pgfplotsretval-1}
    \foreach \i in {0,...,\pdlastpair}{
      \pgfplotstablegetelem{\i}{claim}\of#1 \let\cl\pgfplotsretval
      \pgfplotstablegetelem{\i}{bond}\of#1 \let\bo\pgfplotsretval
      \pgfplotstablegetelem{\i}{writes}\of#1 \let\ch\pgfplotsretval
      \pgfmathtruncatemacro{\y}{3-\cl}
      % A claim that is not released, with a bond that is not held.
      \pgfmathtruncatemacro{\nobond}{and(\cl!=2,\bo!=1)}
      \ifnum\nobond=1 \global\advance\pdpairsout 1 \node[unbonded] at (\bo,\y) {};
      \else \global\advance\pdpairsin 1 \node[occurs] at (\bo,\y) {};\fi
      % A ring: writes were admitted under such a claim.
      \ifnum\ch>0
        \ifnum\bo=1 \draw[written] (\bo,\y) circle[radius=4.6pt];
        \else \draw[written unbonded] (\bo,\y) circle[radius=4.6pt];\fi
      \fi
    }
    \foreach \name [count=\row from 0] in {live,expired,released,abandoned}{
      \node[pd label,text=pdinkmuted,anchor=east] at (-.6,3-\row) {\name};
    }
  }
  \pdpairspanel{\pdpairsreal}
  \foreach \name [count=\x from 0] in {none,held,refunded,slashed}{
    \node[pd label,text=pdinkmuted,anchor=base] at (\x,3.85) {\name};
  }
  \node[pd label,anchor=base east] at (-.6,3.85) {the claim};
  \node[pd label,anchor=base] at (1.5,4.85) {its bond};
  \node[pd title,anchor=north west,align=left] at (3.75,3.4) {the kernel's harbor};
  \node[pd label,anchor=north west,align=left] at (3.75,2.5)
    {\the\pdpairsin\ combinations, each\\in a tinted cell};
  \begin{scope}[yshift=-2.95cm]
    \pdpairspanel{\pdpairsseparate}
    \node[pd title,anchor=north west,align=left] at (3.75,3.4) {claims and funds\\decided separately};
    \node[pd label,anchor=north west,align=left] at (3.75,1.6)
      {\the\pdpairsout\ more, each a claim\\with no held bond};
  \end{scope}
\end{tikzpicture}
\caption{\textbf{In the kernel's harbor a claim and its bond occur in
four combinations.} A row is the claim's status and a column is the status
of its bond. A mark means that some state of the exhaustive test contained
that combination, and a ring means that writes were admitted under such a
claim. The tinted cells are the combinations the reference model allows.
In the lower table, claims and funds are decided by two components that
are each correct, and the red marks are claims with no held bond. The red
ring is writes admitted after the bond was refunded. [internal,
\texttt{harbor stakes -{}-each}]}
\label{fig:bonded-pairs}
\end{figure}


The rows need a word each. A claim is \emph{live} while its lease runs.
Once the lease has ended it has \emph{run out}, and it stays so until its
holder releases it or another grant is given for its path. In the second
case it is \emph{left}, as Wren's claim is after step 4 of
Listing~\ref{alg:bonded-stake-run-left}: nobody released it, and nobody
now can. In the
kernel's harbor a claim and its bond stand together in four ways. The
other harbor comes to six more, and each of them is a claim that was not
released, with a bond that is not held. The red ring is the second break,
step 5 of Listing~\ref{alg:bonded-stake-run-after}: changes admitted under
a claim that had been released. The lower table
shows the first wrong move from each state and no more, because the walk
does not follow a harbor past a move the rules refuse.
Exercise~\ref{ex:bonded-stake-cells} asks which rule keeps each empty cell
of the upper table empty.

\paragraph{Any numbers.}
The walk stops at 2 workers and 4 units. Kani lifts the bound for one
move, as it did for the funds. The property is \texttt{backed} in
Listing~\ref{alg:bonded-stake-kani-step}: a grant that has not been
released has its own stake, and the stake's bond is held. The harness
starts from any path and any funds for which that is true, with every
number in them chosen by the solver, and makes one move.

\pdlistingroom{27}
\pdexcerpt[caption={The property, and the step of the induction as a Kani
harness.},label={alg:bonded-stake-kani-step}]{ch07b-kani-step}

The moves are the ones that touch a grant, or a bond behind one. For a
claim, the harness takes every grant that Chapter~1's rules could give:
the newest grant again if it was not released, or a grant with the next
token. Which of the two a claim gets is decided by code that handles text,
and that code is not under Kani. The walk does run it.

\pdlistingroom{17}
\pdrun[caption={Kani~0.68 on the core's rules for a claim with a bond, and
on three builds with a guard removed.},label={alg:bonded-stake-kani}]{kani-stakes}

The last of the four verdicts is the step. With the base case, that a
path with no grant has no claim to back, it gives the lifecycle's safety
property for any number of moves, on the one thing assumed of Chapter~1's
rules: a claim gives the newest grant again or one with the next token. The first verdict says that a change is
admitted only under a grant with its own stake and a bond held. The third
says that a grant with a new token gets a bond with a new id. The second
is a control. It puts the step's question to a claim that gives the grant
and posts nothing, and that harness fails, as it must.

\paragraph{Each guard removed.}
This build added three guards, and the kernel has a build with each one
taken out. The walk is run on each
(Listing~\ref{alg:bonded-stake-guards}).

\pdlistingroom{10}
\pdrun[caption={The walk again, on three builds of the core with one guard
removed from each.},label={alg:bonded-stake-guards}]{stakes-guards}

Each build fails in a column of its own. Without \texttt{loose}, claims go
bare in 17,120 states, as they did in the harbor made of two parts. Yet no
change is made under one of them, because \texttt{at\_stake} is still there
and asks for the bond. Without \texttt{at\_stake} no claim goes bare, and
7,032 changes are made after a release. The two rows add up to that
harbor's 25,272 wrong moves, and it had neither guard.

The first row is the one to study. Without \texttt{own}, a new claim takes
whatever stake the path's last grant left behind. Kani's verdicts on the
three builds are the lower half of Listing~\ref{alg:bonded-stake-kani}.
That build passes the step. Every claim in it does have a bond held behind
it. What Heron's claim lacks is a bond of its own. It rides on the bond
that Wren left, and when Heron releases, the kernel pays Wren's bond back
to the wallet with nobody having judged it. The walk sees it, as 648
states in which a bond is shared. The harness for the step cannot, because
\texttt{backed} asks whether a bond is there and not whose it is. The last
section met the same shape in its funds
(Listing~\ref{alg:bonded-build-paid-twice}), where every unit was kept and
one bond was paid twice. A property that holds says what
it mentions and no more. Exercise~\ref{ex:bonded-stake-borrowed} follows
the units through.
\pdscope{The walk covers 2 workers and 2 paths at a price of 2 units,
within the bounds that Listing~\ref{alg:bonded-stake-states} prints, and it holds
the harbor to rules written by the same hand as the core. For a harbor that goes wrong it
counts the first wrong move from each state. The harnesses cover one move
on one path, on the functions of the new module, and say nothing of which
grant a claim is given.}

\begin{pdboundary}
\setlength{\parskip}{4pt plus 1pt}
Nobody may judge. A bond whose claim was left stays held. This build has no
move that settles it, and nobody who could make one. The oracles of
\S\ref{sec:pricing:settlement} are not built.

Nobody is asked who they are, as before. The names in the runs say who
would ask. The core takes a worker's name from the kernel on trust, as
Chapter~1's does, and a refund names nobody.

The price is one number for the whole harbor. \S\ref{sec:pricing} wants it
to depend on the work and on who is doing it.

Every bond comes out of one wallet, the project's. The chapter's bond is
posted by the principal that answers for the work, and here two workers
draw on the same funds.

The harbor is held in memory, like the keeper of the last section.
Chapter~1's kernel writes a grant to its database in the step that decides
it. This harbor's grants and bonds are not in that database yet.
\end{pdboundary}
